% Options for packages loaded elsewhere
% Options for packages loaded elsewhere
\PassOptionsToPackage{unicode}{hyperref}
\PassOptionsToPackage{hyphens}{url}
\PassOptionsToPackage{dvipsnames,svgnames,x11names}{xcolor}
%
\documentclass[
  english,
  a4paper,
]{article}
\usepackage{xcolor}
\usepackage{amsmath,amssymb}
\setcounter{secnumdepth}{-\maxdimen} % remove section numbering
\usepackage{iftex}
\ifPDFTeX
  \usepackage[T1]{fontenc}
  \usepackage[utf8]{inputenc}
  \usepackage{textcomp} % provide euro and other symbols
\else % if luatex or xetex
  \usepackage{unicode-math} % this also loads fontspec
  \defaultfontfeatures{Scale=MatchLowercase}
  \defaultfontfeatures[\rmfamily]{Ligatures=TeX,Scale=1}
\fi
\usepackage{lmodern}
\ifPDFTeX\else
  % xetex/luatex font selection
\fi
% Use upquote if available, for straight quotes in verbatim environments
\IfFileExists{upquote.sty}{\usepackage{upquote}}{}
\IfFileExists{microtype.sty}{% use microtype if available
  \usepackage[]{microtype}
  \UseMicrotypeSet[protrusion]{basicmath} % disable protrusion for tt fonts
}{}
\makeatletter
\@ifundefined{KOMAClassName}{% if non-KOMA class
  \IfFileExists{parskip.sty}{%
    \usepackage{parskip}
  }{% else
    \setlength{\parindent}{0pt}
    \setlength{\parskip}{6pt plus 2pt minus 1pt}}
}{% if KOMA class
  \KOMAoptions{parskip=half}}
\makeatother
% Make \paragraph and \subparagraph free-standing
\makeatletter
\ifx\paragraph\undefined\else
  \let\oldparagraph\paragraph
  \renewcommand{\paragraph}{
    \@ifstar
      \xxxParagraphStar
      \xxxParagraphNoStar
  }
  \newcommand{\xxxParagraphStar}[1]{\oldparagraph*{#1}\mbox{}}
  \newcommand{\xxxParagraphNoStar}[1]{\oldparagraph{#1}\mbox{}}
\fi
\ifx\subparagraph\undefined\else
  \let\oldsubparagraph\subparagraph
  \renewcommand{\subparagraph}{
    \@ifstar
      \xxxSubParagraphStar
      \xxxSubParagraphNoStar
  }
  \newcommand{\xxxSubParagraphStar}[1]{\oldsubparagraph*{#1}\mbox{}}
  \newcommand{\xxxSubParagraphNoStar}[1]{\oldsubparagraph{#1}\mbox{}}
\fi
\makeatother

\usepackage{color}
\usepackage{fancyvrb}
\newcommand{\VerbBar}{|}
\newcommand{\VERB}{\Verb[commandchars=\\\{\}]}
\DefineVerbatimEnvironment{Highlighting}{Verbatim}{commandchars=\\\{\}}
% Add ',fontsize=\small' for more characters per line
\usepackage{framed}
\definecolor{shadecolor}{RGB}{255,255,255}
\newenvironment{Shaded}{\begin{snugshade}}{\end{snugshade}}
\newcommand{\AlertTok}[1]{\textcolor[rgb]{0.02,0.51,0.79}{\textbf{#1}}}
\newcommand{\AnnotationTok}[1]{\textcolor[rgb]{0.49,0.02,0.79}{\textbf{#1}}}
\newcommand{\AttributeTok}[1]{\textcolor[rgb]{0.00,0.34,0.68}{#1}}
\newcommand{\BaseNTok}[1]{\textcolor[rgb]{0.54,0.53,0.53}{\textbf{#1}}}
\newcommand{\BuiltInTok}[1]{\textcolor[rgb]{0.39,0.29,0.61}{\textbf{#1}}}
\newcommand{\CharTok}[1]{\textcolor[rgb]{0.57,0.30,0.62}{#1}}
\newcommand{\CommentTok}[1]{\textcolor[rgb]{0.54,0.53,0.53}{#1}}
\newcommand{\CommentVarTok}[1]{\textcolor[rgb]{0.00,0.58,1.00}{#1}}
\newcommand{\ConstantTok}[1]{\textcolor[rgb]{0.12,0.11,0.11}{\textbf{#1}}}
\newcommand{\ControlFlowTok}[1]{\textcolor[rgb]{0.12,0.11,0.11}{\textbf{#1}}}
\newcommand{\DataTypeTok}[1]{\textcolor[rgb]{0.73,0.02,0.79}{#1}}
\newcommand{\DecValTok}[1]{\textcolor[rgb]{0.54,0.53,0.53}{\textit{#1}}}
\newcommand{\DocumentationTok}[1]{\textcolor[rgb]{0.02,0.79,0.25}{#1}}
\newcommand{\ErrorTok}[1]{\textcolor[rgb]{0.02,0.51,0.79}{#1}}
\newcommand{\ExtensionTok}[1]{\textcolor[rgb]{0.00,0.58,1.00}{\textbf{#1}}}
\newcommand{\FloatTok}[1]{\textcolor[rgb]{0.12,0.11,0.11}{\textit{#1}}}
\newcommand{\FunctionTok}[1]{\textcolor[rgb]{0.39,0.29,0.61}{#1}}
\newcommand{\ImportTok}[1]{\textcolor[rgb]{1.00,0.33,0.00}{#1}}
\newcommand{\InformationTok}[1]{\textcolor[rgb]{0.02,0.79,0.25}{#1}}
\newcommand{\KeywordTok}[1]{\textcolor[rgb]{0.12,0.11,0.11}{\textbf{#1}}}
\newcommand{\NormalTok}[1]{\textcolor[rgb]{0.12,0.11,0.11}{#1}}
\newcommand{\OperatorTok}[1]{\textcolor[rgb]{0.12,0.11,0.11}{#1}}
\newcommand{\OtherTok}[1]{\textcolor[rgb]{0.00,0.43,0.16}{#1}}
\newcommand{\PreprocessorTok}[1]{\textcolor[rgb]{0.00,0.43,0.16}{#1}}
\newcommand{\RegionMarkerTok}[1]{\textcolor[rgb]{0.79,0.21,0.02}{#1}}
\newcommand{\SpecialCharTok}[1]{\textcolor[rgb]{0.24,0.68,0.91}{#1}}
\newcommand{\SpecialStringTok}[1]{\textcolor[rgb]{1.00,0.33,0.00}{#1}}
\newcommand{\StringTok}[1]{\textcolor[rgb]{0.75,0.01,0.01}{#1}}
\newcommand{\VariableTok}[1]{\textcolor[rgb]{0.00,0.34,0.68}{#1}}
\newcommand{\VerbatimStringTok}[1]{\textcolor[rgb]{0.75,0.01,0.01}{#1}}
\newcommand{\WarningTok}[1]{\textcolor[rgb]{0.46,0.73,0.89}{\textbf{#1}}}

\usepackage{longtable,booktabs,array}
\newcounter{none} % for unnumbered tables
\usepackage{calc} % for calculating minipage widths
% Correct order of tables after \paragraph or \subparagraph
\usepackage{etoolbox}
\makeatletter
\patchcmd\longtable{\par}{\if@noskipsec\mbox{}\fi\par}{}{}
\makeatother
% Allow footnotes in longtable head/foot
\IfFileExists{footnotehyper.sty}{\usepackage{footnotehyper}}{\usepackage{footnote}}
\makesavenoteenv{longtable}
\usepackage{graphicx}
\makeatletter
\newsavebox\pandoc@box
\newcommand*\pandocbounded[1]{% scales image to fit in text height/width
  \sbox\pandoc@box{#1}%
  \Gscale@div\@tempa{\textheight}{\dimexpr\ht\pandoc@box+\dp\pandoc@box\relax}%
  \Gscale@div\@tempb{\linewidth}{\wd\pandoc@box}%
  \ifdim\@tempb\p@<\@tempa\p@\let\@tempa\@tempb\fi% select the smaller of both
  \ifdim\@tempa\p@<\p@\scalebox{\@tempa}{\usebox\pandoc@box}%
  \else\usebox{\pandoc@box}%
  \fi%
}
% Set default figure placement to htbp
\def\fps@figure{htbp}
\makeatother
\usepackage{svg}



\ifLuaTeX
\usepackage[bidi=basic,shorthands=off]{babel}
\else
\usepackage[bidi=default,shorthands=off]{babel}
\fi
\ifLuaTeX
  \usepackage{selnolig} % disable illegal ligatures
\fi


\setlength{\emergencystretch}{3em} % prevent overfull lines

\providecommand{\tightlist}{%
  \setlength{\itemsep}{0pt}\setlength{\parskip}{0pt}}



 


\usepackage{fancyhdr}
\usepackage{graphicx}
% Search for images relative to project root at depth 0 (root), 1 (content/), and 2 (content/sub/)
\graphicspath{{assets/images/}{../assets/images/}{../../assets/images/}}
\usepackage{xcolor}
\usepackage{fontspec}
\usepackage{lucide-icons}
\setmainfont{P052}
\usepackage[a4paper, top=3cm, bottom=2.5cm, left=2.5cm, right=2.5cm, headheight=60pt, headsep=8pt, footskip=1.5cm]{geometry}
\newlength{\handoutlogooffset}
\AtBeginDocument{%
  \newlength{\handoutcapht}%
  \settoheight{\handoutcapht}{\small A}%
  \setlength{\handoutlogooffset}{\dimexpr1.2cm-\handoutcapht\relax}%
}
% style for Argdown codeblocks
\usepackage[framemethod=TikZ]{mdframed}
\usepackage{xcolor}
\mdfsetup{skipabove=\topskip,skipbelow=\topskip}
\mdfdefinestyle{codedefault}{%
backgroundcolor=black!4,
roundcorner=2pt,
linecolor=black!30,
linewidth=0.5pt,
leftmargin=10, rightmargin=10,
rightline=true,
innerleftmargin=10,innerrightmargin=10}
% redefining the Shaded-Environment of the pandoc-template to allow customised
% style-settings
% Must run at \begin{document} because pandoc defines Shaded after include-in-header
\AtBeginDocument{%
  \ifdefined\Shaded
    \renewenvironment{Shaded}{\centering \begin{mdframed}[style=codedefault]}{\end{mdframed}}%
  \else
    \newenvironment{Shaded}{\centering \begin{mdframed}[style=codedefault]}{\end{mdframed}}%
  \fi
}

% Adjust fontsize for code-blocks if necessary (by redefining the
% Highlighting-Environment of the pandoc-template)
%\DefineVerbatimEnvironment{Highlighting}{Verbatim}{commandchars=\\\{\},fontsize=\small}
\makeatletter
\@ifpackageloaded{float}{}{\usepackage{float}}
\floatstyle{plain}
\@ifundefined{c@chapter}{\newfloat{argdown}{h}{loargdown}}{\newfloat{argdown}{h}{loargdown}[chapter]}
\floatname{argdown}{Argdown Snippet}
\newcommand*\listofargdowns{\listof{argdown}{List of Argdown Snippets}}
\makeatother
\makeatletter
\@ifpackageloaded{float}{}{\usepackage{float}}
\floatstyle{plain}
\@ifundefined{c@chapter}{\newfloat{argdownmap}{h}{loargdownmap}}{\newfloat{argdownmap}{h}{loargdownmap}[chapter]}
\floatname{argdownmap}{Argument Map}
\newcommand*\listofargdownmaps{\listof{argdownmap}{List of Argument Maps}}
\makeatother
\makeatletter
\@ifpackageloaded{tcolorbox}{}{\usepackage[skins,breakable]{tcolorbox}}
\@ifpackageloaded{fontawesome5}{}{\usepackage{fontawesome5}}
\definecolor{quarto-callout-color}{HTML}{909090}
\definecolor{quarto-callout-note-color}{HTML}{0758E5}
\definecolor{quarto-callout-important-color}{HTML}{CC1914}
\definecolor{quarto-callout-warning-color}{HTML}{EB9113}
\definecolor{quarto-callout-tip-color}{HTML}{00A047}
\definecolor{quarto-callout-caution-color}{HTML}{FC5300}
\definecolor{quarto-callout-color-frame}{HTML}{acacac}
\definecolor{quarto-callout-note-color-frame}{HTML}{4582ec}
\definecolor{quarto-callout-important-color-frame}{HTML}{d9534f}
\definecolor{quarto-callout-warning-color-frame}{HTML}{f0ad4e}
\definecolor{quarto-callout-tip-color-frame}{HTML}{02b875}
\definecolor{quarto-callout-caution-color-frame}{HTML}{fd7e14}
\makeatother
\makeatletter
\@ifpackageloaded{caption}{}{\usepackage{caption}}
\AtBeginDocument{%
\ifdefined\contentsname
  \renewcommand*\contentsname{Table of contents}
\else
  \newcommand\contentsname{Table of contents}
\fi
\ifdefined\listfigurename
  \renewcommand*\listfigurename{List of Figures}
\else
  \newcommand\listfigurename{List of Figures}
\fi
\ifdefined\listtablename
  \renewcommand*\listtablename{List of Tables}
\else
  \newcommand\listtablename{List of Tables}
\fi
\ifdefined\figurename
  \renewcommand*\figurename{Figure}
\else
  \newcommand\figurename{Figure}
\fi
\ifdefined\tablename
  \renewcommand*\tablename{Table}
\else
  \newcommand\tablename{Table}
\fi
}
\@ifpackageloaded{float}{}{\usepackage{float}}
\floatstyle{ruled}
\@ifundefined{c@chapter}{\newfloat{codelisting}{h}{lop}}{\newfloat{codelisting}{h}{lop}[chapter]}
\floatname{codelisting}{Listing}
\newcommand*\listoflistings{\listof{codelisting}{List of Listings}}
\makeatother
\makeatletter
\makeatother
\makeatletter
\@ifpackageloaded{caption}{}{\usepackage{caption}}
\@ifpackageloaded{subcaption}{}{\usepackage{subcaption}}
\makeatother
\usepackage{bookmark}
\IfFileExists{xurl.sty}{\usepackage{xurl}}{} % add URL line breaks if available
\urlstyle{same}
\hypersetup{
  pdftitle={Schlussmuster \& Formalisierung},
  pdfauthor={Sebastian Cacean},
  pdflang={en},
  colorlinks=true,
  linkcolor={blue},
  filecolor={Maroon},
  citecolor={Blue},
  urlcolor={Blue},
  pdfcreator={LaTeX via pandoc}}


\title{Schlussmuster \& Formalisierung}
\usepackage{etoolbox}
\makeatletter
\providecommand{\subtitle}[1]{% add subtitle to \maketitle
  \apptocmd{\@title}{\par {\large #1 \par}}{}{}
}
\makeatother
\subtitle{Handout zur Übersicht}
\author{Sebastian Cacean}
\date{2026-08-01}
\begin{document}
% ---------------------------------------------------------------
% Handout style: colored footer on all pages, meta header on p.1
% ---------------------------------------------------------------
\definecolor{handoutline}{HTML}{b5cbd6}

% --- Base page style: only footer, no header rule ---------
\pagestyle{fancy}
\fancyhf{}
\renewcommand{\headrulewidth}{0pt}
\renewcommand{\footrulewidth}{0pt}

\fancyfoot[C]{%
  \parbox{\linewidth}{\centering
    % \textcolor{handoutline}{\rule{\linewidth}{0.4pt}}\\[3pt]%
    \textcolor{handoutline}{\small \thepage}%
  }%
}

% --- First-page style: meta-info left, logo right ---------
\fancypagestyle{firstpage}{%
  \fancyhf{}%
  \renewcommand{\footrulewidth}{0pt}%
  \fancyfoot[C]{%
    \parbox{\linewidth}{\centering
      % \textcolor{handoutline}{\rule{\linewidth}{0.4pt}}\\[3pt]%
      \textcolor{handoutline}{\small \thepage}%
    }%
  }%
}

% --- Preserve default before-body behaviour ---------------
\renewcommand{\maketitle}{%
  \begin{center}%
    {\LARGE\bfseries Schlussmuster \& Formalisierung\par}%
  \end{center}%
  \parbox{\linewidth}{\centering
   \small Sebastian Cacean | sebastian.cacean@hhu.de | 2026-08-01 %
  }%
  % \textcolor{handoutline}{\rule{\linewidth}{0.4pt}}\\[3pt]%
  \vspace{1.5em}%
}
\maketitle
\thispagestyle{firstpage}


\subsection{Karte}\label{karte}

\pandocbounded{\includesvg[keepaspectratio]{index_files/mediabag/svg-xml-base64,PD94b.svg}}

\subsection{Wichtige Definitionen \&
Erläuterungen}\label{wichtige-definitionen-erluxe4uterungen}

\begin{itemize}
\tightlist
\item
  \textbf{Deduktiv gültige Argumente}: Ein Argument ist genau dann
  deduktiv gültig, wenn gilt: Sind die Prämissen wahr, kann die
  Konklusion nicht falsch sein.
\item
  \textbf{Schlüssige Argumente} (stichhaltige Argumente): Ein Argument
  schlüssig genau dann, wenn es deduktiv gültig ist und seine Prämissen
  wahr sind.\footnote{very tedious\ldots{}}
\item
  \textbf{Gültige Schlüsse der Aussagenlogik}: Gültigkeit ergibt sich
  durch die Bedeutung der aussagenlogischen Satzoperatoren
  (\emph{Wenn\ldots, dann\ldots{}}, \emph{nicht}, \emph{oder} und
  \emph{und}).
\item
  Sätze der Form \emph{``Wenn p, dann q.''} (Subjunktionen) drücken
  notwendige und hinreichende Bedingungen aus.

  \begin{itemize}
  \tightlist
  \item
    p ist hinreichend für q und q ist notwendig für p.
  \end{itemize}
\end{itemize}

\subsection{Satzoperatoren der
Aussagenlogik}\label{satzoperatoren-der-aussagenlogik}

{\def\LTcaptype{none} % do not increment counter
\begin{longtable}[]{@{}
  >{\raggedright\arraybackslash}p{(\linewidth - 4\tabcolsep) * \real{0.2222}}
  >{\raggedright\arraybackslash}p{(\linewidth - 4\tabcolsep) * \real{0.3810}}
  >{\raggedright\arraybackslash}p{(\linewidth - 4\tabcolsep) * \real{0.3968}}@{}}
\toprule\noalign{}
\begin{minipage}[b]{\linewidth}\raggedright
Satzoperator
\end{minipage} & \begin{minipage}[b]{\linewidth}\raggedright
Satzform (schematisch)
\end{minipage} & \begin{minipage}[b]{\linewidth}\raggedright
Aussagenlogische Formel
\end{minipage} \\
\midrule\noalign{}
\endhead
\bottomrule\noalign{}
\endlastfoot
Subjunktion (materiales Konditional) &
\texttt{\textquotesingle{}Wenn\ p,\ dann\ q.\textquotesingle{}} &
`\(p\rightarrow q\)' oder '\(p\supset q\)' \\
Bisubjunktion (Bikonditional) &
\texttt{\textquotesingle{}p\ genau\ dann,\ wenn\ q.\textquotesingle{}} &
`\(p\leftrightarrow q\)' oder '\(p\equiv q\)' \\
Disjunktion (Adjunktion) &
\texttt{\textquotesingle{}p\ oder\ q.\textquotesingle{}} &
'\(p\vee q\)' \\
Konjunktion & \texttt{\textquotesingle{}p\ und\ q.\textquotesingle{}} &
`\(p\wedge q\)' oder `\(p.q\)' oder '\(p\&q\)' \\
Negation &
\texttt{\textquotesingle{}Es\ ist\ nicht\ der\ Fall,\ dass\ p.\textquotesingle{}}
& `\(\neg p\)' oder auch '\(\sim p\)' \\
\end{longtable}
}

\subsection{Schlussmuster}\label{schlussmuster}

\subsubsection{Konjunktionseinführung}\label{konjunktionseinfuxfchrung}

\begin{Shaded}
\begin{Highlighting}[]
\AlertTok{\textless{}Argumentschema Konjunktionseinführung\textgreater{}}

\ErrorTok{(1)}\NormalTok{ p.}
\ErrorTok{(2)}\NormalTok{ q.}
\WarningTok{{-}{-} mit Konjunktionseinführung aus 1,2 {-}{-}}
\ErrorTok{(3)}\FloatTok{ *Also*:}\NormalTok{ p und q.}
\end{Highlighting}
\end{Shaded}

\subsubsection{Konjunktionsbeseitigung}\label{konjunktionsbeseitigung}

\begin{Shaded}
\begin{Highlighting}[]
\AlertTok{\textless{}Argumentschema Konjunktionsbeseitigung\textgreater{}}

\ErrorTok{(1)}\NormalTok{ p und q.}
\WarningTok{{-}{-} mit Konjunktionsbeseitigung aus 1 {-}{-}}
\ErrorTok{(3)}\FloatTok{ *Also*:}\NormalTok{ p.}
\end{Highlighting}
\end{Shaded}

\subsubsection{Disjunktionseinführung}\label{disjunktionseinfuxfchrung}

\begin{Shaded}
\begin{Highlighting}[]
\AlertTok{\textless{}Argumentschema Disjunktionseinführung\textgreater{}}

\ErrorTok{(1)}\NormalTok{ p.}
\WarningTok{{-}{-} mit Disjunktionseinführung aus 1,2 {-}{-}}
\ErrorTok{(3)}\FloatTok{ *Also*:}\NormalTok{ p oder q.}
\end{Highlighting}
\end{Shaded}

\subsubsection{Disjunktiver Syllogismus / Adjunktiver
Syllogismus}\label{disjunktiver-syllogismus-adjunktiver-syllogismus}

\begin{Shaded}
\begin{Highlighting}[]
\AlertTok{\textless{}Argumentschema Disjunktiver Syllogismus\textgreater{}}

\ErrorTok{(1)}\NormalTok{ p oder q.}
\ErrorTok{(2)}\NormalTok{ nicht p.}
\WarningTok{{-}{-} aus 1,2 mit disj. Syll. {-}{-}}
\ErrorTok{(3)}\FloatTok{ *Also*:}\NormalTok{ q.}
\end{Highlighting}
\end{Shaded}

Variante mit mehr als 2 Alternativen:

\begin{Shaded}
\begin{Highlighting}[]
\AlertTok{\textless{}Argumentschema Disjunktiver Syllogismus\textgreater{}}

\ErrorTok{(1)}\NormalTok{ p oder r oder ... oder s oder q.}
\ErrorTok{(2)}\NormalTok{ nicht p.}
\ErrorTok{(3)}\NormalTok{ nicht r}
\NormalTok{    ...}
\NormalTok{(n) nicht s}
\WarningTok{{-}{-} aus 1 bis n mit disj. Syll. {-}{-}}
\ErrorTok{(3)}\FloatTok{ *Also*:}\NormalTok{ q.}
\end{Highlighting}
\end{Shaded}

\subsubsection{Kontravalenz}\label{kontravalenz}

\begin{Shaded}
\begin{Highlighting}[]
\AlertTok{\textless{}Argumentschema Kontravalenz\textgreater{}}

\ErrorTok{(1)}\NormalTok{ Entweder p oder q.}
\ErrorTok{(2)}\NormalTok{ p.}
\WarningTok{{-}{-} aus 1,2 mit Kontravalenz{-}{-}}
\ErrorTok{(3)}\FloatTok{ *Also*:}\NormalTok{ q ist falsch.}
\end{Highlighting}
\end{Shaded}

\subsubsection{Prinzip der doppelten
Negation}\label{prinzip-der-doppelten-negation}

\begin{Shaded}
\begin{Highlighting}[]
\AlertTok{\textless{}Argumentschema: Prinzip der doppelten Negation\textgreater{}}

\ErrorTok{(1)}\NormalTok{ Es ist nicht der Fall, dass es nicht der Fall ist, dass p.}
\WarningTok{{-}{-} mit Prinzip der doppelten Negation aus 1 {-}{-}}
\ErrorTok{(2)}\FloatTok{ *Also*:}\NormalTok{ p.}
\end{Highlighting}
\end{Shaded}

\begin{Shaded}
\begin{Highlighting}[]
\AlertTok{\textless{}Argumentschema: Prinzip der doppelten Negation\textgreater{}}

\ErrorTok{(1)}\NormalTok{ p.}
\WarningTok{{-}{-} mit Prinzip der doppelten Negation aus 1 {-}{-}}
\ErrorTok{(2)}\FloatTok{ *Also*:}\NormalTok{ Es ist nicht der Fall, dass es nicht der Fall ist, dass p.}
\end{Highlighting}
\end{Shaded}

\subsubsection{Modus Ponens}\label{modus-ponens}

\begin{Shaded}
\begin{Highlighting}[]
\AlertTok{\textless{}Argumentschema Modus Ponens\textgreater{}}

\ErrorTok{(1)}\NormalTok{ Wenn p, dann q.}
\ErrorTok{(2)}\NormalTok{ p.}
\WarningTok{{-}{-} mit Modus Ponens aus 1,2 {-}{-}}
\ErrorTok{(3)}\FloatTok{ *Also*:}\NormalTok{ q.}
\end{Highlighting}
\end{Shaded}

\subsubsection{Modus Tollens}\label{modus-tollens}

\begin{Shaded}
\begin{Highlighting}[]
\AlertTok{\textless{}Argumentschema Modus Tollens\textgreater{}}

\ErrorTok{(1)}\NormalTok{ Wenn p, dann q.}
\ErrorTok{(2)}\NormalTok{ Es ist nicht der Fall, dass q.}
\WarningTok{{-}{-} mit Modus Tollens aus 1,2 {-}{-}}
\ErrorTok{(3)}\FloatTok{ *Also*:}\NormalTok{ Es ist nicht der Fall, dass p.}
\end{Highlighting}
\end{Shaded}

\subsubsection{Klassisches Dilemma}\label{klassisches-dilemma}

\begin{Shaded}
\begin{Highlighting}[]
\AlertTok{\textless{}Argumentschema Klassisches Dilemma\textgreater{}}

\ErrorTok{(1)}\NormalTok{ Wenn p, dann r.}
\ErrorTok{(2)}\NormalTok{ Wenn nicht p, dann r.}
\WarningTok{{-}{-} aus 1,2, mit klass. Dilemma {-}{-}}
\ErrorTok{(3)}\NormalTok{ Also: r.}
\end{Highlighting}
\end{Shaded}

\subsubsection{Allgemeines Dilemma /
Fallunterscheidung}\label{allgemeines-dilemma-fallunterscheidung}

\begin{Shaded}
\begin{Highlighting}[]
\AlertTok{\textless{}Argumentschema Allgemeines Dilemma\textgreater{}}

\ErrorTok{(1)}\NormalTok{ p oder q.}
\ErrorTok{(2)}\NormalTok{ Wenn p, dann r.}
\ErrorTok{(3)}\NormalTok{ Wenn q, dann r.}
\WarningTok{{-}{-} aus 1{-}3 mit allg. Dilemma {-}{-}}
\ErrorTok{(4)}\FloatTok{ *Also*:}\NormalTok{ r.}
\end{Highlighting}
\end{Shaded}

\begin{Shaded}
\begin{Highlighting}[]
\AlertTok{\textless{}Argumentschema Allgemeines Dilemma\textgreater{}}

\ErrorTok{(1)}\NormalTok{ p oder q oder ... oder ...}
\ErrorTok{(2)}\NormalTok{ Wenn p, dann r.}
\ErrorTok{(3)}\NormalTok{ Wenn q, dann r.}
\ErrorTok{(4)}\NormalTok{ Wenn ..., dann r.}
\ErrorTok{(5)}\NormalTok{ ...}
\WarningTok{{-}{-} aus 1{-}... mit allg. Dilemma {-}{-}}
\ErrorTok{(6)}\FloatTok{ *Also*:}\NormalTok{ r.}
\end{Highlighting}
\end{Shaded}

\subsubsection{Kontraposition}\label{kontraposition}

\begin{Shaded}
\begin{Highlighting}[]
\AlertTok{\textless{}Argumentschema Kontraposition\textgreater{}}

\ErrorTok{(1)}\NormalTok{ Wenn p, dann q.}
\WarningTok{{-}{-} aus 1 mit Kontraposition{-}{-}}
\ErrorTok{(2)}\FloatTok{ *Also*:}\NormalTok{ Wenn q falsch ist, dann ist p falsch.}
\end{Highlighting}
\end{Shaded}

\subsubsection{De Morgan'sche Regeln}\label{de-morgansche-regeln}

Die folgenden Schlussmuster nutzen die folgenden Äquivalenzen aus:

\begin{itemize}
\tightlist
\item
  Konjunktionen der Form \texttt{¬(p∧q)} sind äquivalent zu
  Disjunktionen der Form \texttt{¬p∨¬q}.
\item
  Disjunktionen der Form \texttt{¬(p∨q)} sind äquivalent zu
  Konjunktionen der Form \texttt{¬p∧¬q}.
\end{itemize}

\begin{Shaded}
\begin{Highlighting}[]
\ErrorTok{(1)}\NormalTok{ Es ist nicht der Fall, dass p und q.}
  \CommentTok{// ¬(p∧q)}
\WarningTok{{-}{-} mit De Morgan\textquotesingle{}scher Regel aus 1 {-}{-}}
\ErrorTok{(2)}\FloatTok{ *Also*:}\NormalTok{ Es ist nicht der Fall, dass p oder es ist nicht der Fall, dass q.}
  \CommentTok{// ¬p∨¬q}
\end{Highlighting}
\end{Shaded}

\begin{Shaded}
\begin{Highlighting}[]
\ErrorTok{(1)}\NormalTok{ Es ist nicht der Fall, dass p oder es ist nicht der Fall, dass q.}
  \CommentTok{// ¬p∨¬q}
\WarningTok{{-}{-} mit De Morgan\textquotesingle{}scher Regel aus 1 {-}{-}}
\ErrorTok{(2)}\FloatTok{ *Also*:}\NormalTok{ Es ist nicht der Fall, dass p und q.}
  \CommentTok{// ¬(p∧q)}
\end{Highlighting}
\end{Shaded}

\begin{Shaded}
\begin{Highlighting}[]
\ErrorTok{(1)}\NormalTok{ Es ist nicht der Fall, dass p oder q.}
  \CommentTok{// ¬(p∨q)}
\WarningTok{{-}{-} mit De Morgan\textquotesingle{}scher Regel aus 1 {-}{-}}
\ErrorTok{(2)}\FloatTok{ *Also*:}\NormalTok{ Es ist nicht der Fall, dass p und es ist nicht der Fall, dass q.}
  \CommentTok{// ¬p∧¬q}
\end{Highlighting}
\end{Shaded}

\begin{Shaded}
\begin{Highlighting}[]
\ErrorTok{(1)}\NormalTok{ Es ist nicht der Fall, dass p und es ist nicht der Fall, dass q.}
  \CommentTok{// ¬p∧¬q}
\WarningTok{{-}{-} mit De Morgan\textquotesingle{}scher Regel aus 1 {-}{-}}
\ErrorTok{(2)}\FloatTok{ *Also*:}\NormalTok{ Es ist nicht der Fall, dass p oder q.}
  \CommentTok{// ¬(p∨q)}
\end{Highlighting}
\end{Shaded}

\subsubsection{Transitivität
(Kettenregel)}\label{transitivituxe4t-kettenregel}

\begin{Shaded}
\begin{Highlighting}[]
\AlertTok{\textless{}Argumentschema Transitivität\textgreater{}}

\ErrorTok{(1)}\NormalTok{ Wenn p, dann q.}
\ErrorTok{(2)}\NormalTok{ Wenn q, dann r.}
\WarningTok{{-}{-} aus 1,2 mit Transitivität {-}{-}}
\ErrorTok{(3)}\FloatTok{ *Also*:}\NormalTok{ Wenn p, dann r.}
\end{Highlighting}
\end{Shaded}

\subsubsection{Annahmen um des Arguments
willen}\label{annahmen-um-des-arguments-willen}

\begin{tcolorbox}[enhanced jigsaw, arc=.35mm, bottomrule=.15mm, bottomtitle=1mm, breakable, colback=white, colbacktitle=quarto-callout-note-color!10!white, colframe=quarto-callout-note-color-frame, coltitle=black, left=2mm, leftrule=.75mm, opacityback=0, opacitybacktitle=0.6, rightrule=.15mm, title=\textcolor{quarto-callout-note-color}{\faInfo}\hspace{0.5em}{Annahmen um des Arguments willen}, titlerule=0mm, toprule=.15mm, toptitle=1mm]

Sätze, die im Laufe einer Begründung hypothetisch unterstellt werden,
von deren Wahrheit aber nicht die Wahrheit der Konklusion abhängt,
heißen ``Annahmen um des Arguments willen''

\end{tcolorbox}

\begin{tcolorbox}[enhanced jigsaw, arc=.35mm, bottomrule=.15mm, bottomtitle=1mm, breakable, colback=white, colbacktitle=quarto-callout-important-color!10!white, colframe=quarto-callout-important-color-frame, coltitle=black, left=2mm, leftrule=.75mm, opacityback=0, opacitybacktitle=0.6, rightrule=.15mm, title=\textcolor{quarto-callout-important-color}{\faExclamation}\hspace{0.5em}{Important}, titlerule=0mm, toprule=.15mm, toptitle=1mm]

\begin{itemize}
\tightlist
\item
  Annahmen um des Arguments willen sind \textbf{keine Prämissen}! Sie
  müssen als Annahmen um des Arguments willen gekennzeichnet werden.
\item
  Aus Annahmen um des Arguments willen können -- wie aus Prämissen und
  Zwischenkonklusionen -- unter Anwendung der Schlussregeln
  Zwischenkonklusionen gezogen werden.
\item
  Alle Schlussschritte, in denen Annahmen um des Arguments willen oder
  aber aus diesen bereits gefolgerte Zwischenkonklusionen verwendet
  werden, müssen mit der Nummer der jeweiligen Annahme um des Arguments
  willen gekennzeichnet werden.
\end{itemize}

\end{tcolorbox}

\paragraph{Annahmebeseitigung}\label{annahmebeseitigung}

\begin{Shaded}
\begin{Highlighting}[]
\ErrorTok{(1)}\FloatTok{ *Annahme um des Arguments willen:* }\NormalTok{p}
\NormalTok{   ...}
\WarningTok{{-}{-} folgt mit ... aus ... (hängt ab von 1) {-}{-}}
\NormalTok{(n) q}
\WarningTok{{-}{-} folgt aus 1, n mit Annahmebeseitigung {-}{-}}
\NormalTok{(n+1) Wenn p, dann q.}
\end{Highlighting}
\end{Shaded}

\paragraph{Annahmewiderlegung (reductio ad
absurdum)}\label{annahmewiderlegung-reductio-ad-absurdum}

\begin{Shaded}
\begin{Highlighting}[]
\ErrorTok{(1)}\FloatTok{ *Annahme um des Arguments willen:* }\NormalTok{p}
\NormalTok{   ...}
\WarningTok{{-}{-} folgt mit ... aus ... (hängt ab von 1) {-}{-}}
\NormalTok{(n) q und es ist nicht der Fall, dass q.}
\WarningTok{{-}{-} folgt aus 1,n mit Annahmewiderlegung {-}{-}}
\NormalTok{(n+1) Es ist nicht der Fall, dass p.}
\end{Highlighting}
\end{Shaded}

\subsubsection{Allspezialisierung}\label{allspezialisierung}

\begin{Shaded}
\begin{Highlighting}[]
\AlertTok{\textless{}Argumentschema Allspezialisierung\textgreater{}}

\ErrorTok{(1)}\NormalTok{ Für alle x gilt: x ist F.}
\WarningTok{{-}{-} mit Allspez. aus 1 {-}{-}}
\ErrorTok{(2)}\NormalTok{ Also: Für den Einzelgegenstand a gilt F.}
\end{Highlighting}
\end{Shaded}

Oder:

\begin{Shaded}
\begin{Highlighting}[]
\AlertTok{\textless{}Argumentschema Allspezialisierung\textgreater{}}

\ErrorTok{(1)}\NormalTok{ Für alle x gilt: Wenn x F ist, ist x auch G.}
\WarningTok{{-}{-} mit Allspez. aus 1 {-}{-}}
\ErrorTok{(2)}\FloatTok{ *Also*:}\NormalTok{ Für den Einzelgegenstand a gilt: Ist a F,}
\NormalTok{    ist a auch G.}
\end{Highlighting}
\end{Shaded}

Schematisch angedeutet in einer ganz allgemeinen Form:

\begin{Shaded}
\begin{Highlighting}[]
\AlertTok{\textless{}Argumentschema Allspezialisierung\textgreater{}}

\ErrorTok{(1)}\NormalTok{ Für alle x gilt: ... x ... .}
\WarningTok{{-}{-} mit Allspez. aus 1 {-}{-}}
\ErrorTok{(2)}\FloatTok{ *Also*:}\NormalTok{ Für den Einzelgegenstand a gilt: ... a ... .}
\end{Highlighting}
\end{Shaded}

\subsubsection{Analogieargumente}\label{analogieargumente}

\begin{Shaded}
\begin{Highlighting}[]
\ErrorTok{(1)}\NormalTok{ Gegenstände X verhalten sich in Bezug auf den Aspekt Z* genauso wie}
\NormalTok{    Gegenstände Y in Bezug auf den Aspekt Z.}
\ErrorTok{(2)}\NormalTok{ Auf Gegenstände Y trifft hinsichtlich Z S zu.}
\ErrorTok{(3)}\NormalTok{ Der Sachverhalt S* hinsichtlich Z* entspricht dem Sachverhalt S }
\NormalTok{    hinsichtlich Z.}
\ErrorTok{(4)}\FloatTok{ *Analogieprinzip:* }\NormalTok{Gilt}
\NormalTok{      (i) Gegenstände X verhalten sich in Bezug auf den Aspekt Z* genauso wie}
\NormalTok{      Gegenstände Y in Bezug auf den Aspekt Z}
\NormalTok{      (ii) Auf Gegenstände Y trifft hinsichtlich Z S zu}
\NormalTok{      (iii) Der Sachverhalt S* hinsichtlich Z* entspricht dem Sachverhalt S }
\NormalTok{      hinsichtlich Z}
\NormalTok{    dann gilt auch: Auf Gegenstände X trifft S* zu.}
\WarningTok{{-}{-} aus 1{-}4 mit Allspezialisierung, KE und MP {-}{-}}
\ErrorTok{(5)}\FloatTok{ *Also:* }\NormalTok{Auf Gegenstände X trifft S* zu.}
\end{Highlighting}
\end{Shaded}

Für eine Instantiierung dieses Schemas (also der Formulierung eines
konkreten Analogieargument) müssen Sie für die Prämissen 1-3 und die
Konklusion konkrete Inhalte angeben. Das allgemeine Prinzip (Prämisse 4)
bleibt so stehen (und wird über die Allspezialisierung auf die konkrete
Analogie angewendet).

\subsubsection{Argumentschema: Schluss auf die beste
Erklärung}\label{argumentschema-schluss-auf-die-beste-erkluxe4rung}

\begin{Shaded}
\begin{Highlighting}[]
\ErrorTok{(1)}\NormalTok{ p}
\ErrorTok{(2)}\NormalTok{ Die Annahme, dass q, ist die beste Erklärung für p.}
\WarningTok{{-}{-} aus 1,2 mit dem Schluss auf die beste Erklärung {-}{-}}
\ErrorTok{(3)}\NormalTok{ q}
\end{Highlighting}
\end{Shaded}





\end{document}
